\title{Large Language Models Can Automatically Engineer Features for Few-Shot Tabular Learning}

\begin{abstract}
Large Language Models (LLMs) have demonstrated extraordinary capabilities in handling complex and novel reasoning tasks, suggesting considerable promise for the domain of tabular learning—a foundation for numerous practical applications. This work introduces a new in-context learning paradigm, \textsf{FeatLLM}, in which LLMs serve as feature engineers to construct input datasets optimized for predictive accuracy on tabular data. The features generated via this framework enable the estimation of class probabilities using a straightforward downstream machine learning method such as linear regression, supporting highly effective few-shot learning. Notably, during inference, the \textsf{FeatLLM} approach utilizes only the discovered features together with the simple predictive model, removing reliance on real-time queries to the LLM for each input. This reduces inference overhead compared to prior LLM-driven methods, requires only API access to the LLM, and overcomes constraints imposed by prompt length limitations. Experiments across a variety of tabular benchmarks from multiple domains show that \textsf{FeatLLM} constructs superior predictive rules, outperforming competing approaches like TabLLM and STUNT by an average margin of 10\%.
\end{abstract}

\section{Introduction}
Large Language Models (LLMs) have recently exhibited remarkable proficiency in generating accurate outputs for a broad range of unseen problems, even in the absence of explicit fine-tuning for specific tasks~\cite{creswell2022selection,imani2023mathprompter,taylor2022galactica}. Owing to their training on massive textual datasets, LLMs accumulate rich prior knowledge that can stand in for domain expertise in various fields~\cite{Genomes2021,singhal2023large,wilcox2003role}, thereby facilitating automation in multiple applications, from conversational analysis~\cite{finch2023leveraging} to automated code repair~\cite{fan2023automated}. By embedding task instructions and a handful of instances within prompts, LLMs are capable of discerning context, identifying key features and salient examples~\cite{brown2020language,zhao2023survey}. This underscores LLMs’ analytical competence and hints at their suitability for automated feature engineering.

The prior knowledge and reasoning abilities inherent to LLMs have been increasingly harnessed within tabular learning contexts. For example, leveraging information such as the correlation between age and disease risk can bolster performance in medical prediction scenarios. Recent research has implemented natural language definitions of tasks and features, converted datasets into serialized formats, and then processed them via LLMs~\cite{dinh2022lift,wang2023anypredict}. This is typically coupled with methods such as in-context learning~\cite{nam2023semi} or parameter-efficient fine-tuning~\cite{dinh2022lift,hegselmann2023tabllm}. The exploitation of LLMs’ intrinsic knowledge has demonstrated notable gains over traditional tabular methods, particularly in settings with limited available examples~\cite{hegselmann2023tabllm}.

Existing methods for applying LLMs to tabular data encounter several obstacles. For end-to-end prediction tasks, each data instance requires at least one forward pass through the LLM, resulting in considerable computational overhead. Additionally, achieving strong predictive performance typically necessitates fine-tuning the LLMs, which is impractical for many state-of-the-art models that lack public training infrastructure or provide access only via restricted inference APIs~\cite{anil2023palm,openai2023gpt4}. Another complication is the limited effectiveness of most LLM-based techniques when faced with long prompts~\cite{liu2023lost}; as the dimensionality of tabular data grows and serialized input sequences lengthen, the methods often become infeasible or experience substantial drops in accuracy. Collectively, these issues impede the deployment of LLM-powered tabular learning in practical scenarios.

In contrast, our method reframes the application of LLMs from direct prediction towards enhanced feature construction. Rather than relying on LLMs solely for end-to-end output, we leverage them as advanced feature engineering tools, capitalizing on their embedded knowledge to facilitate more efficient use of information. To this end, we introduce a novel in-context learning strategy, \textsf{FeatLLM}, designed to distill the “criteria” LLMs utilize in making predictions. As an example, in a disease prediction task, an LLM might articulate explicit rules specifying which conditions of the input features lead to a positive diagnosis. These learned criteria, in the form of generated rules, are then used to synthesize new features that replace original attributes for downstream modeling. This approach reassigns the LLM’s role to feature creation, allowing downstream prediction models to be simpler, which subsequently reduces inference time. Through in-context learning, we are able to derive informative features solely by augmenting prompts with illustrative examples, circumventing the need for additional model training. Furthermore, by drawing upon ensemble learning strategies~\cite{breiman2001random}, such as feature bagging, we are able to address scalability issues that arise from overly lengthy prompts.

\textsf{FeatLLM} organizes the process of prompt-based feature engineering into two key stages: first, it focuses on comprehending the underlying problem and deducing how features are linked to the prediction targets; second, leveraging insights from this understanding along with a few provided examples, it develops critical decision rules capable of classifying data points. This approach guides LLMs to disregard irrelevant features during rule formulation. The generated rules are subsequently implemented on the dataset (interpreted as programmatic expressions), transforming each data instance into a binary variable reflecting rule adherence. A straightforward model is then trained using these binary features to predict class membership probabilities. To enhance model reliability, an ensemble method is adopted—this involves repeated feature generation paired with feature bagging, which curtails issues stemming from large prompt sizes by fine-tuning the number of features and data points selected during bagging. \textsf{FeatLLM} does not depend on conventional training procedures and incurs minimal inference costs, facilitating its deployment for a range of practical applications.

For empirical assessment, \textsf{FeatLLM} is tested on 13 tabular datasets under limited data conditions, demonstrating consistently strong and stable results. We find that LLMs are adept at drawing on both domain-specific background and dataset-derived information, thereby identifying effective features. In our comparisons, this framework surpasses existing few-shot learning methods across multiple benchmarks. The source code is accessible via an anonymized GitHub repository at~\texttt{https://github.com/Sungwon-Han/FeatLLM}.

\section{Related Work}

\subsection{Few-Shot Learning with Tabular Data}  
Few-shot learning has made it possible to derive generalizable feature representations from datasets, even when labeled data is scarce~\cite{chen2018closer}. Originally developed for image-based problems, the methods have proven effective across an expanding range of domains, such as textual, audio, and most recently, tabular data formats~\cite{majumder2022few,nam2023stunt,schick2021exploiting}. However, relying on a limited number of labeled instances can increase susceptibility to spurious correlations~\cite{bartlett2021deep}. To address this challenge, recent research has incorporated additional information sources and domain-specific prior knowledge, thereby enforcing beneficial inductive biases in the modeling process. Approaches include leveraging vast pools of unlabeled tabular datasets augmented to support semi-supervised training~\cite{ucar2021subtab,yoon2020vime}, or achieving favorable model initialization through unsupervised pre-training not specific to downstream tasks~\cite{somepalli2021saint}. Unsupervised strategies often rely on masking or corrupting data segments and training models with reconstruction-based objectives~\cite{arik2021tabnet,majmundar2022met}, or exploiting contrastive learning through data augmentation~\cite{bahri2022scarf,chen2020simple}. Nevertheless, the intrinsic diversity of tabular data impedes the identification of universally effective augmentations, with these methods showing limited efficacy under few-shot conditions~\cite{nam2023stunt}.

More recent approaches train models across a wide spectrum of tabular datasets—extending well beyond the intended target tasks—and subsequently transfer this knowledge to the task of interest~\cite{zhang2023generative,levin2022transfer,zhu2023xtab}. As an illustration, TransTab~\cite{wang2022transtab} utilizes transformer architectures to learn semantic associations by incorporating column names, in addition to cell values, thus enriching the representation of table structure. The amassed cross-table knowledge enables models to achieve few-shot or even zero-shot inference on previously unseen tabular problems. Similarly, TabPFN~\cite{hollmann2022tabpfn} constructs synthetic training data using mixtures of distributions to pre-train a Bayesian neural network. Post pre-training, the model is able to make predictions for both training and testing samples in the target task without requiring additional adaptation. The transfer of prior knowledge across a collection of datasets has allowed these methods to outperform classical tree-based models and proved especially advantageous in low-data scenarios.

\subsection{Language-Interfaced Tabular Learning}  
Incorporating prior knowledge via large language models (LLMs)~\cite{anil2023palm,openai2023gpt4} presents another compelling strategy. Benefiting from training on massive and heterogeneous text datasets, LLMs exhibit impressive generalization to novel tasks and display versatility across domains~\cite{brown2020language}. Recent research efforts have explored ways to transfer the knowledge encapsulated in LLMs for tabular learning applications. This generally involves transforming tabular data into text sequences and presenting them as part of the LLM’s input prompts~\cite{dinh2022lift,hegselmann2023tabllm,wang2023anypredict}. Such prompts may include table content, feature annotations, and task descriptions, enabling the model to infer relevant relationships needed for inference. Further progress has involved adapting LLMs specifically to tabular problems through parameter-efficient finetuning techniques like LoRA~\cite{hu2021lora} or IA3~\cite{liu2022few}, as well as harnessing in-context learning—providing the LLM with few-shot exemplars within the input prompt to guide prediction, without additional weight updates~\cite{dinh2022lift,hegselmann2023tabllm,nam2023semi,slack2023tablet}.

While the previously described approaches have shown success in enhancing few-shot learning on tabular datasets, their reliance on routing inference through the LLM for each prediction introduces significant obstacles to industrial implementation. This work seeks to move beyond the typical black-box paradigm, in which predictions for every input sample are generated end-to-end by an LLM. Instead, it proposes utilizing the LLM to deduce explicit rules or conditions that are indicative of each class label. \textsf{FeatLLM} converts these discovered rules into programmatic code, which is then used to engineer new features. This enables downstream predictions to be made independently of the LLM, capitalizing on in-context learning for rule extraction by exploiting prior knowledge alongside a limited number of labeled examples.

\section{Methods}
The core objective of \textsf{FeatLLM} is to identify and extract "rules"—definable sets of conditions corresponding to each possible output class—from a small collection of labeled samples, as illustrated in Figure~\ref{fig:main_model}. Rules identified via the LLM are then translated into binary programmatic features, where each feature represents whether a given rule is satisfied by a sample. These binary features serve as the input for a weighted dot product (with non-negative, trainable parameters), generating a score that reflects the class probability for the sample. Through training, the model not only determines the significance of each feature but also how these features collectively contribute to the probability of each class (see Section~\ref{Sec:training}). The method further increases robustness by employing bagging—repeating the process multiple times—and combining the resulting predictions using an ensemble strategy. The overall problem setup and step-by-step methodology are detailed in the subsequent sections.

\vspace*{-0.12in}
\paragraph{Problem formulation.} We begin with a tabular dataset comprised of $N$ labeled instances, denoted as $\mathcal{D} = \{(\mathbf{x}^i, \mathbf{y}^i)\}_{i=1}^{N}$. Each data point $\mathbf{x}^i$ is a $d$-dimensional vector corresponding to the dataset's $d$ input attributes, which are associated with natural language names $F = \{f_j\}_{j=1}^{d}$ and accompanied by separately supplied definitions. Within the classification context, $\mathbf{y}^i$ represents the label and belongs to a predetermined collection of categories. For $k$-shot learning scenarios, only $k$ instances, where $k < N$, are randomly chosen to serve as the training set for the model. 

\subsection{Prompt Design for Extracting Rules}
\label{Sec:prompt_design}
To steer the LLM toward rule extraction that mirrors an expert's logical reasoning, we construct prompts that encourage it to emulate a knowledgeable human's strategy for interpreting tabular data. The prompt contains three primary sections, as depicted in Fig.~\ref{fig:prompt_for_rule} and detailed in Appendix~\ref{sec:example_rule}:

\vspace*{-0.12in}
\paragraph{Basic information description.} This segment delivers key contextual details needed for addressing the task (highlighted in orange in Fig.~\ref{fig:prompt_for_rule}). It summarizes the problem, specifies the target labels, delineates the feature descriptions, and includes illustrative sample cases. The core question framing the task (e.g., ``Does this patient's myocardial infarction complications data indicate chronic heart failure? Yes or no?") is presented, with further examples available in Appendix~\ref{sec:task_desc}. Feature explanations specify the data type and, where relevant, include the definition; for categorical attributes, example category values are listed. As representative demonstrations, several training samples are converted into serialized text format alongside their correct labels, adopting a serialization style consistent with prior literature~\cite{dinh2022lift,hegselmann2023tabllm}:

\begin{align}
&\text{Serialize}(\mathbf{x}^i,\mathbf{y}^i, F) = \nonumber \\ 
&``\ f_1\ \text{is}\ \mathbf{x}^i_1.\ ... \ f_d\  \text{is}\  \mathbf{x}^i_d.\ \ \text{Answer:}\  \mathbf{y}^i\ ”
\end{align}

\vspace*{-0.12in}
\paragraph{Reasoning instruction.} Our objective is to improve the reasoning capabilities of the LLM by supplying it with tailored reasoning instructions (refer to blue-highlighted text in Fig.~\ref{fig:prompt_for_rule}). These reasoning instructions commence with a prompt akin to that used in the chain-of-thought framework~\cite{wei2022chain}. The subsequent inference process for the LLM is delineated into two distinct phases. During the initial phase, the LLM is prompted to deduce causal links or tendencies between features and the task specification, operating independently from exemplar demonstrations. Instead, the model utilizes general domain knowledge and commonsense reasoning, maximizing the extraction of relevant prior information. In the second phase, the LLM incorporates both the earlier inferred information and the provided demonstrations to formulate class-specific rules. This bifurcated approach is intended to mitigate the risk of the model overfitting to spurious relationships present in uninformative features and supports a focus on attributes with substantial significance. To maintain a manageable and effective set of rules—thus avoiding the twin pitfalls of underfitting due to too few rules or excessively lengthy prompts—we stipulate that the LLM extract a fixed set of rules per class (specifically, 10 rules)\footnote{See Section~\ref{sec:analysis} for further discussion regarding rule quantity.}. Furthermore, when generating these rules, the LLM is prompted to reference each feature's type as indicated in the basic information description, minimizing the potential for type-related inconsistencies.

\vspace*{-0.12in}
\paragraph{Response instruction.} To ensure that the derived rules are both interpretable and easily usable, we instruct the LLM regarding the organization and formatting of its outputs (see yellow-highlighted text in Fig.~\ref{fig:prompt_for_rule}). The prompt provides explicit instructions concerning the formatting requirements for responses associated with each answer category (i.e., the answer format), as well as the expected syntactic structure of each individual rule (i.e., the rule format). Such detailed directives empower the LLM to adapt its response formatting according to the nature of each feature. In the absence of further specification, the LLM is also able to construct composite rules via logical operators such as AND and OR. For an illustrative result, refer to Appendix~\ref{sec:outcome-rule}.

\subsection{Inferring Class Likelihood via Rules}
\label{Sec:training}

\paragraph{Parsing rules for feature generation.} In this stage, the previously generated rules are employed to construct new binary features. For every class, these features signify whether a sample meets the relevant rules (assigning either ‘0’ or ‘1’), which can then serve as predictors for the sample’s class membership. Due to the fact that these rules are derived from natural language responses generated by an LLM, it is necessary to develop an automated parsing method for feature creation. Although we enforce structured formatting for outputs during rule generation to assist parsing, outputs from the LLM often contain syntactic inconsistencies or unexpected variations~\cite{kovavc2023large}. For example, the LLM might replace symbols such as ``$>$" or ``$<$" with their verbal equivalents ``greater than" or ``less than", complicating the parsing process.

To mitigate these difficulties, rather than devising sophisticated parsing algorithms, we opt to use the LLM itself. The LLM’s strength in comprehending semantics regardless of syntax and its capacity to synthesize code from instructions (since it has been trained on extensive code corpora) make it suitable for this task. We formulate prompts that specify the function name, input/output requirements, and the set of rules, which are then provided to the LLM. Examples of these prompts and their resultant outputs are depicted in Figures~\ref{fig:prompt_for_parsing} and~\ref{sec:example_parsing}-\ref{sec:example_parse_outcome} in the Appendix. The Python \texttt{exec()} function is subsequently utilized to run the generated code and effect data transformation.

\vspace*{-0.12in}
\paragraph{Inferring class likelihood.} With these newly created binary rule-based features for each class, an initial approach to estimate how likely a sample belongs to a given class is to tally the count of satisfied rules per class (i.e., summing the binary features assigned to each class). Nevertheless, because different rules contribute differently to predictive performance, it is essential to assign weights reflecting their importance based on training data. We do this by fitting a standard machine learning (ML) model to each class's binary features. One straightforward strategy is to use a linear model without an intercept term. Formally, the binary feature vector for sample $\mathbf{x}^i$ and class $k$ is denoted $\mathbf{z}_k^i \in \{0, 1\}^R$, with $R$ representing the number of rules per class and $c$ the number of total classes. The class probabilities $\mathbf{p}^i$ are then calculated as:

\begin{align}
\text{logit}_k^i &= \max(\mathbf{w}_k, 0) \cdot \mathbf{z}_k^i, \nonumber \\
\mathbf{p}^i &= \text{Softmax}({[\text{logit}_{1}^i, ..., \text{logit}_{c}^i]}). \label{eq:infer_prob}
\end{align}

In this formulation, $\mathbf{w}_k$ denotes the trainable parameter associated with each feature in the linear model, reflecting the feature’s relative importance. By constraining these weights to the positive domain, we guarantee that no feature produced by the LLM contributes to decreasing a class’s predicted probability during training. The logit output is subsequently mapped to class probabilities via the softmax transformation. The parameter $\mathbf{w}_k$ is optimized by minimizing the cross-entropy loss over a small number of labeled training examples.

\vspace*{-0.12in}
\paragraph{Ensembling with bagging.} To further strengthen prediction reliability, we repeat the model-building process several times, resulting in a set of diverse inference models whose collective output forms the basis of the final classification via ensemble averaging. Specifically, for $T$ iterations, we gather the set $\{\mathbf{w}[t]\}_{t=1}^T$ of learned weight vectors, and generate the ensemble prediction by averaging the individual class probability vectors $\mathbf{p}^i$ produced by each $\mathbf{w}[t]$, as formalized in Eq.~\ref{eq:infer_prob}.

Ensemble diversity is encouraged by introducing stochasticity into the rule generation phase. This is accomplished by applying a relatively high temperature to the LLM’s sampling process and by randomizing the sequence in which few-shot exemplars are presented—an ordering known to influence LLM outputs~\cite{lu2022fantastically}\footnote{Further analysis of rule diversity appears in Appendix~\ref{sec:diversity}}. To leverage the resultant variability for stronger predictions, we employ a bagging strategy, selecting random subsets of either features or training samples for each individual model, an approach especially advantageous as the dataset grows in size or dimensionality. The ensemble methodology yields two notable benefits: Firstly, misgenerated rules in some ensemble members may be outweighed by accurate rules produced in others, thereby increasing the system’s robustness. This principle resonates with the self-consistency property of LLMs~\cite{wang2022self}, as frequently recurring rules across separate runs have higher reliability. Secondly, the ensemble approach alleviates constraints imposed by the LLM’s prompt length: bagging reduces the input size per model, supporting stable performance even with extensive tabular data. While any given set of rules may overlook certain feature relationships, collectively, the set of weak learners arising from multiple bagged models can mitigate the effects of incomplete coverage in individual runs.

\section{Experiments}
We assess the effectiveness of \textsf{FeatLLM} on a variety of tabular datasets and perform multiple analyses to deepen our understanding of the framework's behavior. Due to page limitations, further experiments—including explorations with alternative LLMs, qualitative investigations, and more—are provided in the Appendix.

\subsection{Performance Evaluation}
\paragraph{Datasets.}
Our experimental study draws upon 13 tabular datasets focused on either binary or multi-class classification. These datasets include: (1) Adult~\cite{asuncion2007uci}, used to predict whether an individual earns above \$50,000 per year; (2) Bank~\cite{moro2014data}, which targets the likelihood that a customer subscribes to a term deposit; (3) Blood~\cite{yeh2009knowledge}, predicting repeat donation behavior among blood donors; (4) Car~\cite{kadra2021well}, assessing car quality; (5) Communities~\cite{misc_communities_and_crime_183}, forecasting regional crime rates; (6) Credit-g~\cite{kadra2021well}, determining creditworthiness of individuals; (7) Diabetes\footnote{\texttt{kaggle.com/datasets/uciml/pima-indians-diabetes-database}}, which aims to detect diabetes in a patient; (8) Heart\footnote{\texttt{kaggle.com/datasets/fedesoriano/heart-failure-prediction}}, focused on predicting coronary artery disease occurrences; and (9) Myocardial~\cite{misc_myocardial_infarction_complications_579}, which identifies individuals suffering from chronic heart failure.

To complement these, we introduce contemporary and synthetic datasets not represented in LLM pre-training corpora: (10) Cultivars, intended to predict soybean cultivar grain yield\footnote{\texttt{archive.ics.uci.edu/dataset/913}}; (11) NHANES, predicting a person's age category from given features\footnote{\texttt{archive.ics.uci.edu/dataset/887}}; (12) Sequence-type, a synthetic classification problem distinguishing sequence types—arithmetic, geometric, Fibonacci, or Collatz; and (13) Solution-mix, tasked with deciding if the concentration in a solution mixture exceeds 0.5. The first two datasets were recently contributed to the UCI Machine Learning Repository (post-September 2023), precluding their presence in the pre-training data of the LLM API (gpt-3.5-turbo-0613) leveraged in this study. The latter two are constructed synthetic datasets, ensuring the LLM API could not have previously encountered them, thereby enabling a robust evaluation.

These datasets encompass a spectrum of sizes and levels of complexity, as summarized in Table~\ref{Tab:basic-desc}. Each dataset comes with descriptive attribute names and meanings. Notably, datasets like `Communities' and `Myocardial' feature in excess of 100 attributes, thus imposing constraints given the limited context window in LLMs.

\vspace*{-0.12in}
\paragraph{Baselines.}
We benchmark our method against nine baseline models. The first group encompasses traditional tabular learning algorithms: (1) Logistic Regression (LogReg), (2) XGBoost~\cite{chen2016xgboost}, and (3) RandomForest~\cite{ho1995random}. The subsequent two approaches, (4) SCARF~\cite{bahri2022scarf} and (5) STUNT~\cite{nam2023stunt}, leverage unlabeled data. It is important to note, however, that obtaining large-scale unlabeled tabular data can be impractical in many deployment contexts. An additional recent approach is (6) TabPFN~\cite{hollmann2022tabpfn}, which utilizes diverse synthetic distributions to pre-train its model. The remaining three baselines involve large language models: (7) In-context learning (In-context)~\cite{wei2022emergent}, (8) TABLET~\cite{slack2023tablet}, and (9) TabLLM~\cite{hegselmann2023tabllm}. In-context learning operates by placing a few labeled training instances within the input prompt, foregoing parameter updates. TABLET augments the prompt with supplementary signals, derived from an external classifier that supplies rule-based and prototype-driven cues to enhance inference reliability. TabLLM applies parameter-efficient methods such as IA3~\cite{liu2022few} to fine-tune the LLM with limited labeled examples.

\vspace*{-0.12in}
\paragraph{Implementation details.}
\textsf{FeatLLM} adopts GPT-3.5 as its underlying LLM, though the framework is compatible with various LLM choices (refer to Appendix~\ref{sec:backbones} for experiments using alternate backbones). For LLM inference, a temperature of 0.5 is used, and the API's top-p parameter remains at its default setting of 1. The ensemble size and the count of extraction rules are configured to 20 and 10, respectively. Additional analysis of hyper-parameter choices is presented in Figure~\ref{fig:hyperparameter2} and discussed in Appendix~\ref{sec:hyper-parameter}. The linear model is optimized using Adam with a learning rate of 0.01, and training proceeds for 200 epochs. To identify the optimal number of training epochs, $k$-fold cross-validation is utilized. In baseline implementations, in-context learning is performed with GPT-3.5, while T0~\cite{sanh2022multitask} is adopted for TabLLM, in line with its original protocol. Notably, TabLLM limits LLM use to those models with available open-source checkpoints; GPT-3.5, therefore, is excluded under this criterion. For traditional machine learning approaches such as LogReg, XGBoost, and RandomForest, hyper-parameters are chosen via a combination of Grid-search and $k$-fold cross-validation. The choice of $k$ (either 2 or 4) ensures that every class is represented in the training folds. Further technical specifics about these implementations can be found in Appendix~\ref{sec:baselines}.

\vspace*{-0.12in}
\paragraph{Results.}
Tables~\ref{tab:main1} and \ref{tab:main2} showcase the comparative results across various datasets. Each experiment is conducted three times, each with different random splits for training and testing sets; we report the mean and standard deviation for the resulting AUC scores. Across these benchmarks, \textsf{FeatLLM} either achieves the highest ranking or secures the second-best performance relative to the other baseline methods. Remarkably, our method matches the results of traditional tabular learning baselines using 32 shots, with only 4 shots required (see Fig.~\ref{fig:compares}). While STUNT and TabPFN delivered notable gains for select datasets, their aggregate advantage over established machine learning methods was limited. Furthermore, LLM-based strategies—such as in-context learning, TABLET, and TabLLM—were unable to handle datasets characterized by high feature dimensionality. Our method’s use of bagging and ensemble strategies enables robust performance, even for datasets with large numbers of features. Additionally, our ensemble and bagging methodology could, in theory, be integrated into other LLM-based frameworks. Nonetheless, this adaptation would necessitate twice as many inferences per data instance, which would significantly elevate computational requirements, rendering such an approach impractical.

\subsection{Analysis \& Discussion}
\label{sec:analysis}
\paragraph{Ablation study.} To investigate the impact of key methodological components, we conducted ablation studies targeting: (1) \textsf{-Tuning}: removing the linear model’s weight tuning and instead calculating logits by summing binary features; (2) \textsf{-Ensemble}: omitting the ensemble strategy; (3) \textsf{-Description}: excluding feature description; and (4) \textsf{-Reasoning}: eliminating the Step-1 reasoning process during rule generation.

Table~\ref{tab:ablation} reports the average AUC differences from these ablations over all datasets. In every scenario, excluding any of these components led to diminished performance. Notably, the contribution of weight tuning amplifies as the number of shots increases, suggesting that having more training examples allows for more accurate rule weighting, thereby making tuning substantially more effective. Conversely, both feature description and reasoning instructions have greater positive impact at lower shot numbers. This finding underscores the importance of effectively leveraging the prior knowledge encoded in LLMs when only limited data is available. Finally, our proposed ensemble technique consistently provides performance enhancements across all experimental configurations.

\vspace*{-0.12in}
\paragraph{Training \& inference time comparison.} We evaluate the training and inference durations across multiple approaches. All experiments utilize the Adult dataset, specifically using a training set composed of 16 shots. While a single A100 GPU is generally used for evaluation, TabLLM is an exception and relies on four A100 GPUs to facilitate model parallelism. Table~\ref{tab:cost} summarizes these time-based comparisons. Importantly, \textsf{FeatLLM} achieves a notably low inference time that is on par with that of standard, non-LLM methods. On the other hand, LLM-centric techniques such as In-context learning, TABLET, and TabLLM incur substantially greater inference times. When it comes to training durations, \textsf{FeatLLM} remains comparable with alternative LLM-based models; in particular, it necessitates only 30 API calls, representing a minor computational cost that can also be distributed across parallel processes.

\vspace*{-0.12in}
\paragraph{Quality of features generated by \textsf{FeatLLM}.}
To assess the effectiveness of rule-based features produced by \textsf{FeatLLM}, we carry out a straightforward experiment to examine their informativeness for real-world tasks. Concretely, we calculate the average AUROC between the novel features and the true labels, and then identify the proportion of these features that achieve an AUROC above 0.5 for each dataset. The outcomes, presented in Table~\ref{tab:quality}, illustrate that most of the generated rules are meaningfully connected to the target variable and contribute valuable information to the primary task (see the “Before tuning” column). Furthermore, when fitting a linear model to the data, rules exhibiting minimal relevance (i.e., those whose estimated weights fall below a preset threshold) are excluded during optimization (refer to the “After tuning” column). The threshold for pruning was set to 0.1, determined in light of the weight distribution.

\vspace*{-0.12in}
\paragraph{Effect of adding spurious correlations.} To evaluate how well different approaches can eliminate noisy spurious correlations under limited labeled data, we performed a targeted study. Specifically, using the Heart dataset as a base, we randomly appended columns drawn from the Adult dataset that have no relevance to the Heart classification objective. We then tracked model performance as the number of irrelevant features was gradually increased. To maintain fairness, every method was allotted the same 8 labeled examples and we did not incorporate any unlabeled data, thereby omitting methods such as SCARF and STUNT from the comparison.

Figure~\ref{fig:spurious} displays the resulting changes in accuracy stemming from the inclusion of spurious features. Among the approaches evaluated, \textsf{FeatLLM} exhibited the smallest decline in accuracy in the face of irrelevant data. This robustness can be attributed to our method’s requirement that rules align with prior domain knowledge concerning feature-task relevance, thereby discouraging the inclusion of rules that depend on unrelated features and supporting consistent performance. Notably, only a minor proportion (1–2\%) of the rules relied on the artificially introduced noisy columns. Nevertheless, we also find that \textsf{FeatLLM} remains vulnerable to spurious correlations, particularly when features from datasets with inherent semantic connections are randomly included (see Appendix~\ref{sec:spurious-appendix}).

\vspace*{-0.12in}
\paragraph{Hyper-parameter analysis}
\textsf{FeatLLM} incorporates two primary hyper-parameters: the ensemble count and the quantity of rules extracted per class in each LLM query. To evaluate the effect of these settings, we performed a series of experiments. Our findings indicate that increasing the ensemble size enhances performance, yet the improvements plateau when the ensemble count exceeds roughly 20 (refer to Fig.~\ref{fig:appendix-hyperparameter1} in Appendix). In terms of rule quantity, although differences were not statistically significant, setting the number to 10 was deemed optimal given the balance between prompt length and predictive performance (see Fig.~\ref{fig:hyperparameter2}). We conjecture that including additional rules expands the input dimensionality, supplying more information but also heightening the risk of overfitting in the low-shot scenario.

\section{Conclusion}
We introduce \textsf{FeatLLM}, a framework that establishes a new benchmark for LLM-based few-shot tabular learning. Rather than simply using LLMs for sample-wise predictions, \textsf{FeatLLM} leverages LLMs as feature engineers to generate predictive criteria. By employing rule extraction and aggregating through bagging, \textsf{FeatLLM} offers extremely low inference overhead, operates solely via API access with no model retraining, and is adaptable to constraints on feature dimensions. The method delivers robust predictive accuracy and provides a highly effective, data-efficient solution for deploying LLMs on real-world tabular datasets.

\textsf{FeatLLM} is specifically tailored for settings with scarce labeled examples. Moving forward, we aim to extend its feature construction techniques to accommodate larger datasets, investigate alternative types of features beyond rule-based approaches, and further integrate interpretability—thereby broadening its application potential across diverse domains.

\section{Broader Impact} The \textsf{FeatLLM} framework is designed to exploit the prior knowledge embedded in LLMs for tabular tasks, enabling improvements surpassing standard supervised approaches under limited data conditions. By enhancing data efficiency, our work helps address the prevalent challenge of overfitting associated with small training sets, enriching learning with pre-existing knowledge. \textsf{FeatLLM} stands to be particularly valuable in applied contexts such as finance or healthcare, where acquiring labeled data is resource-intensive or logistically difficult, and LLMs’ pretraining on relevant domain texts provides a significant advantage. Nevertheless, for more widespread and responsible adoption, it will be crucial to systematically assess and address how LLM-induced biases or inaccuracies, potentially inherited from their pretraining corpora, might influence or override decisions relative to directly observed or collected ground truth labels.

\end{document}